\section{Conclusion}\label{sec:conclusion}

This report treated robot localization and SLAM as dynamic Bayesian networks and the particle filter as approximate
inference on them: the exact filtering recursion with importance sampling in place of the integrals, and for SLAM a
Rao-Blackwellized split into a sampled trajectory and per-landmark Kalman filters, justified by the conditional
independence of the landmarks given the trajectory. On this basis it compared the sampling choices of the filter in
controlled experiments with 20 simulated worlds per configuration.

The number of particles and the use of resampling decided most of the accuracy. The four resampling schemes were equally
accurate in localization, with systematic resampling keeping the most distinct particles and, in FastSLAM with few
particles, giving a better trajectory and map than multinomial resampling. Resampling only when the effective sample
size or the largest weight signals degeneracy was as accurate as resampling at every step while resampling at a third
to a half of the steps with low thresholds. Proposals that use the current measurement kept more particles effective
and helped most when particles are few: the auxiliary particle filter was the most accurate localization filter from
100 particles on, and FastSLAM 2.0 built better maps than FastSLAM 1.0 with 5 and 10 particles. The extended Kalman
particle filter, in the form tested here, did not improve on the motion model. For the localization world, a few hundred
particles gave most of the achievable accuracy.

The study also showed that reused filter code has to be checked before it is compared: the original FastSLAM 2.0 code
did not sample from its proposal, and several other errors would have biased the results. Natural extensions are MCMC
moves after resampling, the linearized optimal proposal for the extended Kalman particle filter, the resampling rules in
SLAM, and real data with unknown correspondences.
